% Lapisan solusi O001 -- Bab 4 (Ruang Hilbert)
% 10 solusi asli terpisah; bukan tulisan John M. Erdman.
% Provenans model: OpenAI Codex gpt-5.6-sol, Ultra, atas arahan pengguna.
% Lisensi komponen solusi: CC BY-SA 4.0.
% Tidak menyiratkan dukungan John M. Erdman atau Portland State University.

\markboth{SOLUSI O001 UNTUK BAB 4: RUANG HILBERT}{SOLUSI O001 UNTUK BAB 4: RUANG HILBERT}
\section*{Solusi O001 untuk Bab 4: Ruang Hilbert}
\addcontentsline{toc}{section}{Solusi O001 untuk Bab 4: Ruang Hilbert}
\begin{center}
\small
Komponen solusi asli terpisah, CC BY-SA 4.0. Ditulis dengan bantuan
OpenAI Codex gpt-5.6-sol, Ultra, atas arahan pengguna. Pernyataan latihan
berasal dari edisi terjemahan karya John M. Erdman; solusi ini bukan tulisan
beliau dan tidak menyiratkan dukungan beliau atau Portland State University.
\end{center}

% O001-SOLUTION-ID: O001-FAOA-2015-CH04-EX-001-SOLUTION
% SOURCE-EXERCISE-ID: FAOA-2015-CH04-NODE-0023
% STATEMENT-TARGET-SHA256: a4f5325f6a535b95e804661fb3f2a11e0b9b8d5a613e45903338613c2f6e4c70
\begin{o001solution}
  {O001-FAOA-2015-CH04-EX-001-SOLUTION}
  {FAOA-2015-CH04-NODE-0023}
  {a4f5325f6a535b95e804661fb3f2a11e0b9b8d5a613e45903338613c2f6e4c70}
\begin{o001statement}
 Apakah jumlah langsung yang didefinisikan di atas (dengan morfisme yang sesuai) merupakan hasil kali dalam kategori $\cat{HSp}$
yang objeknya ruang Hilbert dan morfismenya pemetaan linear terbatas?  Apakah jumlah langsung itu merupakan koproduk?
\end{o001statement}
\begin{o001answer}
Jumlah langsung ortogonal suatu keluarga $(H_i)_{i\in I}$ merupakan hasil kali
sekaligus koproduk dalam $\cat{HSp}$ jika dan hanya jika hanya berhingga banyak
$H_i$ yang tak nol.
\end{o001answer}
\begin{o001proof}
Jika hanya $H_1,\ldots,H_m$ yang tak nol, proyeksi koordinat
$p_i\colon\bigoplus_iH_i\to H_i$
bersifat terbatas. Jika $f_i\colon K\to H_i$ terbatas, maka
\[
 f(x)=(f_1x,\ldots,f_mx)
\]
adalah satu-satunya pemetaan yang memenuhi $p_i f=f_i$, dan
$\norm{f}^2\leq\sum_i\norm{f_i}^2$. Jadi jumlah langsung berhingga adalah
hasil kali. Dengan inklusi koordinat $j_i$, untuk pemetaan terbatas
$g_i\colon H_i\to K$, rumus
$g(x_1,\ldots,x_m)=\sum_i g_i x_i$ memberikan satu-satunya morfisme dengan
$g j_i=g_i$. Jadi objek yang sama juga koproduk.

Sebaliknya, andaikan tak hingga banyak $H_i$ yang tak nol. Pilih indeks
berbeda $i_1,i_2,\ldots$ dan vektor satuan $u_n\in H_{i_n}$. Untuk sifat
hasil kali, ambil $K=\K$, definisikan
$f_{i_n}(\alpha)=\alpha u_n$, dan ambil $f_i=0$ pada indeks lain. Setiap
$f_i$ terbatas. Namun pemetaan yang dipaksakan oleh semua proyeksi harus
mengirim $1$ ke keluarga yang koordinat $i_n$-nya sama dengan~$u_n$ untuk
semua~$n$; keluarga itu tidak dapat dijumlahkuadratkan. Jadi morfisme
pemfaktor tidak ada.

Untuk sifat koproduk, definisikan
$g_{i_n}(x)=\langle x,u_n\rangle$ dan $g_i=0$ pada indeks lain. Setiap
$g_i$ terbatas. Jika morfisme pemfaktor $g$ ada, maka untuk
\[
 x_N=\frac1{\sqrt N}\sum_{n=1}^{N}j_{i_n}u_n
\]
berlaku $\|x_N\|=1$, tetapi $g(x_N)=\sqrt N$. Ini bertentangan dengan
keterbatasan~$g$. Maka kedua sifat universal gagal tepat ketika terdapat tak
hingga banyak summan tak nol.
\end{o001proof}
\end{o001solution}

% O001-SOLUTION-ID: O001-FAOA-2015-CH04-EX-002-SOLUTION
% SOURCE-EXERCISE-ID: FAOA-2015-CH04-NODE-0053
% STATEMENT-TARGET-SHA256: 22c06694e8a8ce55ab43731d60ce3c9656a0b7af895ba5cbb039796129395f13
\begin{o001solution}
  {O001-FAOA-2015-CH04-EX-002-SOLUTION}
  {FAOA-2015-CH04-NODE-0053}
  {22c06694e8a8ce55ab43731d60ce3c9656a0b7af895ba5cbb039796129395f13}
\begin{o001statement}
Suatu fungsi $f\colon \R \sto \C$ berbentuk
   \[ f(t) = a_0 + \sum_{k=1}^n (a_k \cos kt + b_k \sin kt) \]
dengan $a_0, \dots, a_n, b_1, \dots, b_n$ bilangan kompleks disebut
 \index{polinom trigonometri}%
 \index{polinom!trigonometri}%
\df{polinom trigonometri}.
 \begin{enumerate}
  \item Jelaskan mengapa fungsi bernilai kompleks yang $2\pi$-periodik pada garis real $\R$
sering diidentifikasi dengan fungsi bernilai kompleks pada lingkaran satuan~$\T$.
  \item Beberapa penulis menyebut \emph{polinom trigonometri} sebagai fungsi berbentuk
    \[ f(t) = \sum_{k=-n}^n c_k e^{ikt} \]
 dengan $c_{-n}, \dots, c_{-1}, c_0, c_1, \dots, c_n$ bilangan kompleks. Jelaskan dengan cermat mengapa bentuk ini persis
sama dengan definisi sebelumnya.
  \item Jelaskan alasan penggunaan istilah \emph{polinom trigonometri}.
  \item Buktikan bahwa setiap fungsi kontinu yang $2\pi$-periodik pada $\R$ merupakan limit seragam
dari suatu barisan polinom trigonometri.
 \end{enumerate}
\end{o001statement}
\begin{o001answer}
(1) Identifikasi itu melalui $t\mapsto e^{it}$ dan ruang hasil bagi
$\R/(2\pi\Z)\cong\T$. (2) Kedua bentuk terkait oleh
\[
 c_0=a_0,\qquad c_k=\frac{a_k-i b_k}{2},\qquad
 c_{-k}=\frac{a_k+i b_k}{2}\quad(k\geq1).
\]
(3) Bentuk eksponensial merupakan polinom Laurent pada $z=e^{it}$.
(4) Jumlah Fej\'er dari deret Fourier $f$ adalah polinom trigonometri dan
konvergen seragam ke~$f$.
\end{o001answer}
\begin{o001proof}
Pemetaan $q\colon\R\to\T$, $q(t)=e^{it}$, bernilai sama tepat pada kelas
kongruensi modulo $2\pi$. Karena itu setiap fungsi $2\pi$-periodik mempunyai
faktorisasi tunggal melalui $\R/(2\pi\Z)$, yang secara alami diidentifikasi
dengan~$\T$.

Identitas Euler memberikan
\[
 \cos kt=\frac{e^{ikt}+e^{-ikt}}2,
 \qquad
 \sin kt=\frac{e^{ikt}-e^{-ikt}}{2i}.
\]
Substitusi memberi rumus koefisien dalam jawaban. Sebaliknya,
$a_k=c_k+c_{-k}$ dan $b_k=i(c_k-c_{-k})$, sehingga kedua keluarga fungsi
persis sama. Dengan $z=e^{it}$, bentuk itu adalah pembatasan
$\sum_{k=-n}^n c_kz^k$ pada $\abs z=1$; inilah alasan istilah
``polinom trigonometri''.

Tinggal dibuktikan pernyataan aproksimasi. Untuk $f$ kontinu dan
$2\pi$-periodik, tetapkan
\[
 \widehat f(k)=\frac1{2\pi}\int_{-\pi}^{\pi}f(s)e^{-iks}\,ds
\]
dan jumlah Fej\'er
\[
 \sigma_N f(t)=\sum_{k=-N}^{N}\left(1-\frac{\abs k}{N+1}\right)
                 \widehat f(k)e^{ikt}.
\]
Ini polinom trigonometri. Perkalian hingga dan perubahan variabel langsung
memberi
\[
 \sigma_N f(t)=\frac1{2\pi}\int_{-\pi}^{\pi}F_N(u)f(t-u)\,du,
 \qquad
 F_N(u)=\frac1{N+1}
 \left(\frac{\sin((N+1)u/2)}{\sin(u/2)}\right)^2,
\]
dengan nilai kontinu $F_N(0)=N+1$. Kernel ini tak negatif dan
$\frac1{2\pi}\int_{-\pi}^{\pi}F_N=1$. Maka
\[
 \abs{\sigma_N f(t)-f(t)}
 \leq\frac1{2\pi}\int_{-\pi}^{\pi}
 F_N(u)\abs{f(t-u)-f(t)}\,du.
\]
Untuk $\epsilon>0$, kekontinuan seragam $f$ memberi $\delta\in(0,\pi)$
sehingga faktor selisih kurang dari $\epsilon/2$ apabila $\abs u<\delta$.
Pada $\delta\leq\abs u\leq\pi$ berlaku
\[
 F_N(u)\leq\frac1{(N+1)\sin^2(\delta/2)}.
\]
Bagian integral tersebut jadi paling besar
$2\norm f_\infty/((N+1)\sin^2(\delta/2))$, terlepas dari~$t$, dan kurang
dari $\epsilon/2$ untuk $N$ cukup besar. Jadi
$\norm{\sigma_Nf-f}_\infty\to0$.
\end{o001proof}
\end{o001solution}

% O001-SOLUTION-ID: O001-FAOA-2015-CH04-EX-003-SOLUTION
% SOURCE-EXERCISE-ID: FAOA-2015-CH04-NODE-0078
% STATEMENT-TARGET-SHA256: 4f8e6fe2fc3fbbf3c5cd40c97d7a3c2295778e82b32add1bf6cf200560d0a1b4
\begin{o001solution}
  {O001-FAOA-2015-CH04-EX-003-SOLUTION}
  {FAOA-2015-CH04-NODE-0078}
  {4f8e6fe2fc3fbbf3c5cd40c97d7a3c2295778e82b32add1bf6cf200560d0a1b4}
\begin{o001statement}
Pada suatu sore, kawan Anda, Fred R. Dimm, datang membawa masalah. ``Begini,'' katanya,
``Pada $\lfs 2 = \lfs 2(\R,\R)$, fungsional evaluasi $E_0$, yang mengevaluasi setiap anggota $\lfs 2$
di~$0$, jelas merupakan fungsional linear terbatas. Jadi, menurut Teorema Representasi Riesz,
seharusnya ada fungsi $g$ dalam $\lfs 2$ sedemikian sehingga $f(0) = \langle f,g \rangle = \int_{-\infty}^\infty fg$
untuk semua $f$ dalam~$\lfs 2$. Namun, itu adalah sifat `fungsi~$\delta$', yang katanya
tidak ada. Di mana letak kesalahan saya?'' Berilah Freddy nasihat (yang membantu).
\end{o001statement}
\begin{o001answer}
$E_0$ bukan fungsional pada $\lfs2(\R)$: anggota $\lfs2$ adalah kelas
ekuivalensi hampir-di-mana-mana, sehingga nilai di satu titik tidak
terdefinisi. Bahkan pada fungsi kontinu terintegralkan-kuadrat, evaluasi di
nol tidak terbatas terhadap norma $\lfs2$.
\end{o001answer}
\begin{o001proof}
Jika dua fungsi hanya berbeda nilainya di nol, keduanya menentukan anggota
$\lfs2(\R)$ yang sama, tetapi ``evaluasi di nol'' akan memberi nilai berbeda.
Jadi $E_0$ tidak terdefinisi pada ruang itu.

Masalah tersebut tidak hilang dengan memilih fungsi kontinu sebagai wakil.
Untuk $n\geq1$, ambil
\[
 f_n(t)=\max\{1-n\abs t,0\}.
\]
Fungsi ini kontinu, $f_n(0)=1$, dan perhitungan langsung memberi
\[
 \norm{f_n}_2^2
 =2\int_0^{1/n}(1-nt)^2\,dt=\frac{2}{3n}\longrightarrow0.
\]
Jika evaluasi terbatas terhadap norma $\lfs2$, seharusnya
$1=\abs{E_0f_n}\leq\norm{E_0}\norm{f_n}_2\to0$, suatu kontradiksi.
Karena hipotesis Teorema Representasi Riesz gagal, tidak ada vektor
$g\in\lfs2$ yang merepresentasikan evaluasi titik tersebut.
\end{o001proof}
\end{o001solution}

% O001-SOLUTION-ID: O001-FAOA-2015-CH04-EX-004-SOLUTION
% SOURCE-EXERCISE-ID: FAOA-2015-CH04-NODE-0079
% STATEMENT-TARGET-SHA256: 7a531d47180bfa90ec00b09db0665bc412ad4270a85442b46f63ce4efa10ab53
\begin{o001solution}
  {O001-FAOA-2015-CH04-EX-004-SOLUTION}
  {FAOA-2015-CH04-NODE-0079}
  {7a531d47180bfa90ec00b09db0665bc412ad4270a85442b46f63ce4efa10ab53}
\begin{o001statement}
Masih pada sore yang sama, Freddy kembali. Setelah mempertimbangkan
nasihat yang Anda berikan (dalam soal sebelumnya), ia merevisi
pertanyaannya dengan menjadikan domain $E_0$ sebagai himpunan fungsi bernilai real yang terbatas dan
kontinu pada~$\R$. Dengan sabar Anda
menjelaskan kepada Freddy bahwa sekarang setidaknya ada dua kekeliruan dalam caranya
menggunakan `fungsi~$\delta$'. Apakah kedua kekeliruan tersebut?
\end{o001statement}
\begin{o001answer}
Dengan norma supremum, $C_b(\R)$ memang membuat $E_0$ terbatas, tetapi
$C_b(\R)$ bukan ruang Hilbert, sehingga teorema representasi Hilbert--Riesz
tidak berlaku. Jika ia justru bermaksud memakai hasil kali dalam
$\int fg$, hasil kali itu bahkan tidak terdefinisi pada semua $C_b(\R)$;
setelah domain dibatasi agar masuk $L_2$, evaluasi tetap tidak terbatas.
\end{o001answer}
\begin{o001proof}
Pada $C_b(\R)$ dengan norma supremum berlaku
$\abs{f(0)}\leq\norm f_\infty$, jadi evaluasi memang kontinu. Namun norma
supremum tidak berasal dari hasil kali dalam: misalnya fungsi konstan
$f=1$ dan $g$ dengan nilai dalam $[0,1]$, $g(0)=1$, dan tumpuan kompak
dapat dipilih sehingga
\[
 \norm{f+g}_\infty^2+\norm{f-g}_\infty^2=4+1\neq4
 =2\norm f_\infty^2+2\norm g_\infty^2.
\]
Jadi hukum jajargenjang gagal dan Teorema Representasi Riesz untuk ruang
Hilbert tidak dapat diterapkan.

Di sisi lain, rumus $\langle f,g\rangle=\int_\R fg$ tidak menentukan hasil
kali dalam pada seluruh $C_b(\R)$: fungsi konstan $1$ saja tidak berada
dalam $L_2(\R)$. Jika domain diganti dengan $C_b(\R)\cap L_2(\R)$ dan norma
$L_2$, barisan tenda $f_n$ dari solusi sebelumnya masih memenuhi
$f_n(0)=1$ dan $\norm{f_n}_2\to0$. Maka evaluasi tidak terbatas. Evaluasi
titik secara benar direpresentasikan oleh ukuran titik Dirac dalam teori
representasi untuk fungsional pada ruang fungsi kontinu, bukan oleh suatu
fungsi $L_2$ dalam argumen Freddy.
\end{o001proof}
\end{o001solution}

% O001-SOLUTION-ID: O001-FAOA-2015-CH04-EX-005-SOLUTION
% SOURCE-EXERCISE-ID: FAOA-2015-CH04-NODE-0080
% STATEMENT-TARGET-SHA256: 40661e351b7cad48baf7c9d8e4b91c629468788e4e93a6722bd384dc6146a546
\begin{o001solution}
  {O001-FAOA-2015-CH04-EX-005-SOLUTION}
  {FAOA-2015-CH04-NODE-0080}
  {40661e351b7cad48baf7c9d8e4b91c629468788e4e93a6722bd384dc6146a546}
\begin{o001statement}
Mimpi buruk itu berlanjut. Sekarang pukul 23.00 dan Freddy kembali lagi.
Kali ini (setelah, untungnya, menyerah soal fungsi~$\delta$) ia ingin
menerapkan teorema representasi pada fungsional ``evaluasi
pada nol'' pada ruang~$l_2(\Z)$. Dapatkah ia melakukannya? Jelaskan.
\end{o001statement}
\begin{o001answer}
Ya. Evaluasi koordinat nol adalah fungsional terbatas bernorma satu pada
$l_2(\Z)$ dan direpresentasikan oleh vektor basis standar $e^0$:
$E_0(x)=x_0=\langle x,e^0\rangle$.
\end{o001answer}
\begin{o001proof}
Berbeda dari $L_2(\R)$, anggota $l_2(\Z)$ benar-benar merupakan barisan
dengan koordinat yang terdefinisi. Untuk setiap $x=(x_k)_{k\in\Z}$,
\[
 \abs{x_0}^2\leq\sum_{k\in\Z}\abs{x_k}^2=\norm x_2^2,
\]
sehingga $E_0$ linear dan $\norm{E_0}\leq1$. Pada $e^0$ berlaku
$\norm{e^0}=1$ dan $E_0(e^0)=1$, jadi normanya tepat satu. Dengan konvensi
hasil kali dalam edisi ini,
\[
 \langle x,e^0\rangle=\sum_{k\in\Z}x_k\conj{e^0_k}=x_0.
\]
Jadi semua hipotesis teorema representasi terpenuhi dan representernya
adalah~$e^0$.
\end{o001proof}
\end{o001solution}

% O001-SOLUTION-ID: O001-FAOA-2015-CH04-EX-006-SOLUTION
% SOURCE-EXERCISE-ID: FAOA-2015-CH04-NODE-0091
% STATEMENT-TARGET-SHA256: 986b9bb676932d7ef63f7ab92f49a8fa51ab6068b73b4daedb6c0c01a7c2ec38
\begin{o001solution}
  {O001-FAOA-2015-CH04-EX-006-SOLUTION}
  {FAOA-2015-CH04-NODE-0091}
  {986b9bb676932d7ef63f7ab92f49a8fa51ab6068b73b4daedb6c0c01a7c2ec38}
\begin{o001statement}
Jelaskan mengapa topologi pada ruang Hilbert yang dibangkitkan oleh konvergensi lemah jaring sesungguhnya merupakan
topologi lemah dalam pengertian topologis yang lazim.
\end{o001statement}
\begin{o001answer}
Topologi itu tepat sama dengan topologi awal
$\sigma(H,H^*)$: topologi terkecil yang membuat setiap fungsional linear
terbatas kontinu.
\end{o001answer}
\begin{o001proof}
Untuk setiap $y\in H$, definisikan $\phi_y(x)=\langle x,y\rangle$. Teorema
Representasi Riesz menyatakan bahwa keluarga $\{\phi_y:y\in H\}$ adalah
seluruh dual kontinu~$H^*$. Topologi awal yang ditentukan keluarga ini
mempunyai basis lingkungan di $a$ berupa
\[
 U(a;y_1,\ldots,y_m;\epsilon)
 =\{x:\abs{\langle x-a,y_j\rangle}<\epsilon,
             \ 1\leq j\leq m\}.
\]
Menurut sifat umum topologi awal, suatu jaring $(x_\lambda)$ konvergen ke
$a$ dalam topologi tersebut jika dan hanya jika
$\phi_y(x_\lambda)\to\phi_y(a)$ untuk setiap $y\in H$. Kondisi ini persis
$\langle x_\lambda,y\rangle\to\langle a,y\rangle$ untuk setiap~$y$, yaitu
definisi konvergensi lemah dalam soal. Karena konvergensi jaring menentukan
topologi secara tunggal, topologi yang dibangkitkan oleh konvergensi itu
adalah $\sigma(H,H^*)$, topologi lemah dalam arti topologis.
\end{o001proof}
\end{o001solution}

% O001-SOLUTION-ID: O001-FAOA-2015-CH04-EX-007-SOLUTION
% SOURCE-EXERCISE-ID: FAOA-2015-CH04-NODE-0092
% STATEMENT-TARGET-SHA256: 848170d90169079965a64c432bbc9a6b1bfaf0ab895ba7558340a4a60524926c
\begin{o001solution}
  {O001-FAOA-2015-CH04-EX-007-SOLUTION}
  {FAOA-2015-CH04-NODE-0092}
  {848170d90169079965a64c432bbc9a6b1bfaf0ab895ba7558340a4a60524926c}
\begin{o001statement}
Misalkan $H$ ruang Hilbert dan $(x_{{}_{\sst\lambda}})$ jaring dalam $H$.
 \begin{enumerate}
  \item Jika $x_{{}_{\sst\lambda}} \to^s  a$ dalam $H$, maka $x_{{}_{\sst\lambda}} \to^w a$ dalam $H$.
  \item Topologi kuat (norma) pada $H$ lebih kuat (lebih besar) daripada topologi lemah.
  \item Tunjukkan dengan contoh bahwa kebalikan dari (a) tidak berlaku.
  \item Apakah norma pada $H$ kontinu secara lemah?
  \item Jika $x_{{}_{\sst\lambda}} \to^w  a$ dalam $H$ dan jika $\norm{x_{{}_{\sst\lambda}}} \to \norm{a}$, maka
$x_{{}_{\sst\lambda}}\to^s a$ dalam~$H$.
  \item Jika pemetaan linear $T\colon H \sto H$ kontinu (secara kuat), maka pemetaan tersebut kontinu secara lemah.
 \end{enumerate}
\end{o001statement}
\begin{o001answer}
(1), (2), (5), dan (6) benar. Untuk (3), basis standar $(e^n)$ di $l_2$
konvergen lemah ke nol tetapi tidak dalam norma. Untuk (4), norma kontinu
secara lemah tepat ketika $H$ berdimensi hingga.
\end{o001answer}
\begin{o001proof}
(1) Dari ketaksamaan Cauchy--Schwarz,
\[
 \abs{\langle x_\lambda-a,y\rangle}
 \leq\norm{x_\lambda-a}\norm y\longrightarrow0
\]
untuk setiap $y\in H$. (2) Jadi setiap jaring yang konvergen dalam topologi
norma juga konvergen lemah; ekuivalen dengan mengatakan bahwa topologi norma
lebih kuat.

(3) Dalam $l_2$, untuk setiap $y=(y_n)\in l_2$ berlaku
$\langle e^n,y\rangle=\conj{y_n}\to0$, sebab koordinat setiap barisan
$l_2$ menuju nol. Namun $\norm{e^n}=1$, sehingga $e^n$ tidak konvergen dalam
norma ke nol.

(4) Contoh yang sama menunjukkan bahwa, jika $H$ berdimensi tak hingga,
norma tidak kontinu lemah di nol: pilih suatu barisan ortonormal $(e^n)$;
maka $e^n\to^w0$ tetapi $\norm{e^n}=1\not\to0$. Jika $H$ berdimensi hingga,
ambil basis ortonormal $e^1,\ldots,e^m$. Konvergensi lemah berarti konvergensi
setiap koordinat, dan
$\norm x^2=\sum_{j=1}^m\abs{\langle x,e^j\rangle}^2$; jadi topologi lemah
sama dengan topologi norma dan fungsi norma kontinu.

(5) Karena hasil kali dalam kontinu pada setiap koordinat lemah yang tetap,
\[
 \norm{x_\lambda-a}^2
 =\norm{x_\lambda}^2+\norm a^2
  -2\Re\langle x_\lambda,a\rangle\longrightarrow0.
\]
(6) Untuk $y\in H$, pemetaan
$x\mapsto\langle Tx,y\rangle$ adalah fungsional linear terbatas, sebab nilai
mutlaknya paling besar $\norm T\norm y\norm x$. Oleh Riesz, ada $z_y\in H$
sehingga $\langle Tx,y\rangle=\langle x,z_y\rangle$. Jika
$x_\lambda\to^w a$, maka
$\langle Tx_\lambda,y\rangle\to\langle Ta,y\rangle$ untuk setiap~$y$;
jadi $Tx_\lambda\to^wTa$. Kriteria jaring untuk kontinuitas menyelesaikan
pembuktian.
\end{o001proof}
\end{o001solution}

% O001-SOLUTION-ID: O001-FAOA-2015-CH04-EX-008-SOLUTION
% SOURCE-EXERCISE-ID: FAOA-2015-CH04-NODE-0093
% STATEMENT-TARGET-SHA256: 4c8d88aed6251f2083c6029080a78a9b8a175cd792a4d2e114382a46518d4f4f
\begin{o001solution}
  {O001-FAOA-2015-CH04-EX-008-SOLUTION}
  {FAOA-2015-CH04-NODE-0093}
  {4c8d88aed6251f2083c6029080a78a9b8a175cd792a4d2e114382a46518d4f4f}
\begin{o001statement}
Misalkan $H$ ruang Hilbert berdimensi tak hingga. Tinjau subhimpunan-subhimpunan berikut dari~$H$:
 \begin{enumerate}[label=\rm{(\alph*)}]
  \item bola satuan terbuka;
  \item bola satuan tertutup;
  \item sfer satuan;
  \item subruang linear tertutup.
 \end{enumerate}
Tentukan untuk setiap himpunan tersebut apakah masing-masing bersifat: tertutup secara kuat, tertutup secara lemah,
terbuka secara kuat, terbuka secara lemah, kompak secara kuat, kompak secara lemah. (Salah satu bagian
terlalu sulit untuk saat ini: \emph{Teorema Alaoglu}~\futurexref{6.2.9}{C067441} akan menunjukkan bahwa
bola satuan tertutup kompak secara lemah.)
\end{o001statement}
\begin{o001answer}
Dengan urutan sifat ``tertutup kuat, tertutup lemah, terbuka kuat, terbuka
lemah, kompak kuat, kompak lemah'', jawabannya adalah
\[
\begin{array}{c|cccccc}
 &\mathrm{TK}&\mathrm{TL}&\mathrm{BK}&\mathrm{BL}&\mathrm{KK}&\mathrm{KL}\\ \hline
\{\norm x<1\}&\text{tidak}&\text{tidak}&\text{ya}&\text{tidak}&\text{tidak}&\text{tidak}\\
\{\norm x\leq1\}&\text{ya}&\text{ya}&\text{tidak}&\text{tidak}&\text{tidak}&\text{ya}\\
\{\norm x=1\}&\text{ya}&\text{tidak}&\text{tidak}&\text{tidak}&\text{tidak}&\text{tidak}.
\end{array}
\]
Untuk subruang linear tertutup $M$: ia tertutup kuat dan lemah; ia terbuka
kuat maupun lemah tepat jika $M=H$; dan ia kompak kuat maupun lemah tepat
jika $M=\{0\}$.
\end{o001answer}
\begin{o001proof}
Kontinuitas norma memberi keterbukaan bola terbuka serta ketertutupan bola
tertutup dan sfer dalam topologi kuat. Ketiganya tidak mempunyai sifat kuat
lain yang tercantum: bola terbuka tidak tertutup; dua himpunan lainnya tidak
terbuka; dan tidak satu pun kompak kuat. Untuk klaim terakhir, suatu barisan
ortonormal terletak di bola tertutup dan sfer serta mempunyai jarak
$\sqrt2$ antarsuku. Kelipatan tetap, misalnya $e^n/2$, memberi argumen yang
sama di dalam bola terbuka.

Bola tertutup bersifat tertutup lemah karena
\[
 \{x:\norm x\leq1\}
 =\bigcap_{\norm y\leq1}\{x:\abs{\langle x,y\rangle}\leq1\},
\]
suatu irisan himpunan tertutup lemah. Ia juga kompak lemah. Berikut bukti
lokal klaim terakhir. Petakan bola itu ke
$\prod_{y\in H}\{z\in\C:\abs z\leq\norm y\}$ melalui
$x\mapsto(\langle x,y\rangle)_y$. Produk cakram kompak ini kompak menurut
teorema kekompakan hasil kali. Setiap titik tutupan koordinat menentukan
fungsional konjugat-linear $\Phi(y)$ yang memenuhi
$\abs{\Phi(y)}\leq\norm y$ dan mempertahankan penjumlahan serta perkalian
skalar; sifat-sifat itu adalah persamaan tertutup dalam produk. Teorema
Representasi Riesz memberi $x\in H$, $\norm x\leq1$, dengan
$\Phi(y)=\langle x,y\rangle$. Jadi citra bola tertutup sendiri tertutup di
dalam produk kompak dan karenanya kompak; topologi yang diwarisi tepat
topologi lemah.

Setiap lingkungan lemah dasar hanya membatasi berhingga banyak hasil kali
dalam. Karena $H$ berdimensi tak hingga, terdapat vektor satuan $v$ yang
ortogonal terhadap semua vektor penguji itu. Maka lingkungan tersebut
memuat $x+tv$ untuk setiap skalar~$t$. Akibatnya bola terbuka, bola tertutup,
dan sfer tidak terbuka lemah. Barisan ortonormal memenuhi $e^n\to^w0$;
karena itu bola terbuka tidak tertutup lemah (kelipatan
$(1-1/n)e^1$ juga memberi limit batas), dan sfer tidak tertutup lemah.
Dalam ruang Hausdorff, himpunan kompak harus tertutup; jadi bola terbuka dan
sfer tidak kompak lemah.

Terakhir, untuk subruang tertutup $M$, teorema proyeksi memberi
$M=M^{\perp\perp}$. Dengan demikian
\[
 M=\bigcap_{y\in M^\perp}\{x:\langle x,y\rangle=0\},
\]
sehingga $M$ tertutup lemah. Subruang yang mempunyai interior dalam suatu
topologi ruang vektor harus sama dengan seluruh ruang; jadi $M$ terbuka
dalam salah satu topologi tepat jika $M=H$. Jika $M$ memuat $u\neq0$, fungsi
lemah-kontinu $x\mapsto\langle x,u\rangle$ tidak terbatas pada
$\{tu:t\in\K\}\subset M$, sedangkan citra kontinu himpunan kompak harus
kompak dan terbatas. Jadi $M$ tidak kompak lemah, dan terlebih lagi tidak
kompak kuat. Subruang nol adalah kompak dalam kedua topologi.
\end{o001proof}
\end{o001solution}

% O001-SOLUTION-ID: O001-FAOA-2015-CH04-EX-009-SOLUTION
% SOURCE-EXERCISE-ID: FAOA-2015-CH04-NODE-0094
% STATEMENT-TARGET-SHA256: 6ac5283840818281ac10974dbec5da9b992478b82ef530203afa9f84a1cfcb77
\begin{o001solution}
  {O001-FAOA-2015-CH04-EX-009-SOLUTION}
  {FAOA-2015-CH04-NODE-0094}
  {6ac5283840818281ac10974dbec5da9b992478b82ef530203afa9f84a1cfcb77}
\begin{o001statement}
Misalkan $\{e^n\colon n \in \N\}$ basis biasa untuk $l_2$. Untuk
setiap $n \in \N$, misalkan $a_n = \sqrt n\, e^n$. Manakah di antara pernyataan berikut
yang benar?
 \begin{enumerate}
  \item $0$ adalah titik akumulasi lemah dari himpunan
$\{a_n\colon n \in \N\}$.
  \item $0$ adalah titik gugus lemah dari barisan~$(a_n)$.
  \item $0$ adalah limit lemah dari suatu subbarisan dari~$(a_n)$.
 \end{enumerate}
\end{o001statement}
\begin{o001answer}
Pernyataan (1) dan (2) benar, sedangkan (3) salah.
\end{o001answer}
\begin{o001proof}
Ambil lingkungan lemah dasar nol yang ditentukan oleh
$y^1,\ldots,y^m\in l_2$ dan $\epsilon>0$. Tetapkan
$s_n=\sum_{j=1}^m\abs{y^j_n}^2$. Karena $\sum_ns_n<\infty$, untuk setiap
$N$ ada $n\geq N$ dengan $ns_n<\epsilon^2$: jika tidak, deret ekor
$\sum_ns_n$ mendominasi kelipatan deret harmonik. Untuk indeks itu,
\[
 \abs{\langle a_n,y^j\rangle}
 =\sqrt n\,\abs{y^j_n}\leq\sqrt{ns_n}<\epsilon
 \qquad(1\leq j\leq m).
\]
Jadi setiap lingkungan lemah nol memuat $a_n$ untuk indeks yang arbitrer
besar. Karena $0$ bukan salah satu $a_n$, nol adalah titik akumulasi lemah
himpunannya, dan karena kunjungan terjadi setelah setiap indeks, nol juga
titik gugus lemah barisan.

Sekarang ambil sebarang subbarisan $(a_{n_k})$. Jika ia konvergen lemah ke
$x$, maka untuk setiap koordinat tetap $r$, akhirnya $n_k\neq r$, sehingga
$\langle a_{n_k},e^r\rangle=0$ dan $\langle x,e^r\rangle=0$. Kelengkapan
basis memaksa $x=0$. Pilih lagi subsubbarisan dengan
$n_{k_j}\geq2^j$, dan definisikan $y\in l_2$ dengan
$y_{n_{k_j}}=n_{k_j}^{-1/2}$ serta koordinat lain nol. Memang
$\sum_j\abs{y_{n_{k_j}}}^2\leq\sum_j2^{-j}<\infty$, tetapi
\[
 \langle a_{n_{k_j}},y\rangle=1
\]
untuk setiap~$j$. Subsubbarisan itu tidak konvergen lemah ke nol, sehingga
subbarisan semula pun tidak mungkin konvergen lemah. Maka (3) salah.
\end{o001proof}
\end{o001solution}

% O001-SOLUTION-ID: O001-FAOA-2015-CH04-EX-010-SOLUTION
% SOURCE-EXERCISE-ID: FAOA-2015-CH04-NODE-0097
% STATEMENT-TARGET-SHA256: 4d967e20b7456eb8ff5a320ebaa0e09a669d541a90efd55b495fa12f0ecb5237
\begin{o001solution}
  {O001-FAOA-2015-CH04-EX-010-SOLUTION}
  {FAOA-2015-CH04-NODE-0097}
  {4d967e20b7456eb8ff5a320ebaa0e09a669d541a90efd55b495fa12f0ecb5237}
\begin{o001statement}
Dalam definisi sebelumnya, rujukan kepada funktor pelupa sering
dihilangkan dan diagram yang menyertainya sering digambar sebagai berikut:
 \[\xy
     \qtriangle/{>}`{>}`{-->}/[S`F`A;\iota`f`\wt f_\iota]
   \endxy\]
Diagram ini tentu terlihat jauh lebih sederhana. Menurut Anda, mengapa saya memilih versi yang
lebih rumit?
\end{o001statement}
\begin{o001answer}
Versi lengkap menjaga tipe setiap objek dan panah: $S,\abs F,\abs A$ serta
$\iota,f,\abs{\widetilde f_\iota}$ berada di kategori himpunan, sedangkan
$F,A$ dan $\widetilde f_\iota$ berada di kategori konkret. Diagram ringkas
mencampurkan dua kategori dan menyembunyikan peran funktor pelupa.
\end{o001answer}
\begin{o001proof}
Dalam kategori konkret $\cat C$, $F$ dan $A$ adalah objek kategori, bukan
secara harfiah himpunan. Pemetaan $\iota\colon S\to\abs F$ dan
$f\colon S\to\abs A$ hanyalah fungsi antarhimpunan dan pada umumnya bukan
morfisme $\cat C$. Sebaliknya, pemetaan universal
$\widetilde f_\iota\colon F\to A$ adalah morfisme dalam $\cat C$, dan
persamaan yang benar pada himpunan dasar ialah
\[
 \abs{\widetilde f_\iota}\circ\iota=f.
\]
Diagram sederhana menulis $S,F,A$ seolah-olah ketiganya objek satu kategori
dan menulis ketiga panah seolah-olah bertipe sama. Diagram yang lebih rumit
memperlihatkan baik segitiga pada himpunan dasar maupun morfisme kategori
yang dipetakan ke sana oleh funktor pelupa. Karena itu komutativitas dan
sifat universalnya bertipe benar dan tidak bergantung pada penyalahgunaan
notasi yang hanya aman setelah pembaca memahami identifikasi tersebut.
\end{o001proof}
\end{o001solution}
